% ==============================================================
%  Group Decision-Making and Survival Tasks
% ==============================================================

\section{Group Decision-Making and Survival Tasks}
\label{sec:group-decision-making}

The systems reviewed in Section~\ref{sec:ai-roles} place an AI agent inside group conversations, but the question of what it means for a group to decide well belongs to a different field: group decision-making, studied for decades in social and organisational psychology. Understanding this literature is needed both to frame the problem and to recognise the experimental paradigms the discipline has developed to study it.

The starting point is Steiner's \emph{Group Process and Productivity}~\cite{steiner1972group}, which introduces the notion of \emph{process loss}: a group typically produces less than the sum of its parts, because some individual effort is dissipated in coordination, motivational problems, and inefficient communication. Hill's meta-analysis of decades of comparative studies~\cite{hill1982group} confirms that groups do not systematically beat the best individual: they do under specific conditions and on specific kinds of tasks, but not in general. Stasser and Stewart~\cite{stasser1992hidden} add a specific mechanism known as the \emph{hidden profile} paradigm: when critical information is unevenly distributed among members, so that each one has a unique piece and only by pooling the pieces can the group reach the correct solution, discussion tends to revolve around the information already shared rather than surfacing what is unique, and groups end up making poor decisions. This is the structure of the Murder Mystery task used by Dubini in Section~\ref{sec:ai-roles}. The overall picture is that group performance is contingent: it depends on the task, on the structure of the information, and on the discussion process, which leaves room for procedural guidance, human or artificial, to improve it.

To study these phenomena empirically the discipline has developed \emph{survival decision tasks}: scenarios in which a group has to rank a list of items by survival priority, first individually and then as a group, with performance measured as the deviation from a reference ranking produced by an expert. The original task is NASA Moon Survival, developed by Hall in the early 1960s and formalised in the experimental study by Hall and Watson~\cite{hall1970normative}: fifteen items to rank for survival after an emergency landing on the Moon, with an expert ranking supplied by NASA. Hall and Watson report a result that is central to the present discussion: groups receiving explicit consensus-seeking instructions (look for disagreement, avoid majority voting, do not yield just to avoid conflict) produce significantly better decisions than groups without instructions, and beat their most competent member more often (strong synergy). Nemiroff and Pasmore~\cite{nemiroff1975lostatsea} propose a sibling task, Lost at Sea, with the same structure but a maritime scenario. The paradigm is still active in contemporary research: Hamada and colleagues~\cite{hamada2020wisdom} use NASA Moon Survival to compare group deliberation with wisdom-of-crowds aggregation across twenty-five face-to-face groups, and Hémon and colleagues~\cite{hemon2024constructive} use a survival task in videoconference settings to show that structured instructions based on constructive controversy boost synergy in online groups as well.
